%%%%% PROJECT-SPECIFIC MACROS %%%%%
% Define your project-specific terms here
% Import with: %%%%% PROJECT-SPECIFIC MACROS %%%%%
% Define your project-specific terms here
% Import with: %%%%% PROJECT-SPECIFIC MACROS %%%%%
% Define your project-specific terms here
% Import with: %%%%% PROJECT-SPECIFIC MACROS %%%%%
% Define your project-specific terms here
% Import with: \input{commands/macros}
%
% INSTRUCTIONS:
% 1. Use \textsc{} for method/dataset/model names (small caps)
% 2. Add \xspace at the end for automatic spacing
% 3. Group related terms together
% 4. Use consistent naming conventions

\usepackage{xspace}

%%%%%%%%%%%%%%%%%%%%%%%%%%%%%%%%%%%%%%%%%%%%%%%%%%%%%%%%%%%%%%%%%
% METHOD NAMES
%%%%%%%%%%%%%%%%%%%%%%%%%%%%%%%%%%%%%%%%%%%%%%%%%%%%%%%%%%%%%%%%%

% Your method name (replace with actual name)
\newcommand{\END}{\textsc{end}\xspace}          % Elicitation--Novelty Discriminator

% Discriminator quantities
\newcommand{\SG}{\textsc{sg}\xspace}            % Self-Generation utility
\newcommand{\LCG}{\textsc{lcg}\xspace}          % Length-Controlled Gain

% Task / setting names
\newcommand{\zeth}{\textsc{Zeth}\xspace}
\newcommand{\vocabzeth}{\textsc{vocab-zeth}\xspace}
\newcommand{\vocabes}{\textsc{vocab-es}\xspace}
\newcommand{\arithbten}{\textsc{arith-b10}\xspace}
\newcommand{\arithbnine}{\textsc{arith-b9}\xspace}
\newcommand{\arithbeleven}{\textsc{arith-b11}\xspace}

% Model names
\newcommand{\claude}{\textsc{Claude-Sonnet-4.5}\xspace}
\newcommand{\gpt}{\textsc{GPT-4.1}\xspace}
\newcommand{\qwen}{\textsc{Qwen2.5-7B-Instruct}\xspace}

% Abbreviations
\newcommand{\ie}{i.e.,\xspace}
\newcommand{\eg}{e.g.,\xspace}
\newcommand{\cf}{cf.\xspace}

%%%%%%%%%%%%%%%%%%%%%%%%%%%%%%%%%%%%%%%%%%%%%%%%%%%%%%%%%%%%%%%%%
% DATASET NAMES
%%%%%%%%%%%%%%%%%%%%%%%%%%%%%%%%%%%%%%%%%%%%%%%%%%%%%%%%%%%%%%%%%

% Dataset names (use \textsc for consistency)
% Example: \newcommand{\datasetone}{\textsc{Dataset-One}\xspace}
% Example: \newcommand{\datasettwo}{\textsc{Dataset-Two}\xspace}

%%%%%%%%%%%%%%%%%%%%%%%%%%%%%%%%%%%%%%%%%%%%%%%%%%%%%%%%%%%%%%%%%
% MODEL NAMES
%%%%%%%%%%%%%%%%%%%%%%%%%%%%%%%%%%%%%%%%%%%%%%%%%%%%%%%%%%%%%%%%%

% LLM and model names
% Example: \newcommand{\gpt}{\textsc{GPT-4o-mini}\xspace}
% Example: \newcommand{\llama}{\textsc{Llama-3.1-70B-Instruct}\xspace}
% Example: \newcommand{\claude}{\textsc{Claude-3.5-Sonnet}\xspace}

%%%%%%%%%%%%%%%%%%%%%%%%%%%%%%%%%%%%%%%%%%%%%%%%%%%%%%%%%%%%%%%%%
% BASELINE NAMES
%%%%%%%%%%%%%%%%%%%%%%%%%%%%%%%%%%%%%%%%%%%%%%%%%%%%%%%%%%%%%%%%%

% Baseline method names
% Example: \newcommand{\baseline}{\textsc{BaselineMethod}\xspace}
% Example: \newcommand{\fewshot}{\textsc{Few-shot}\xspace}

%%%%%%%%%%%%%%%%%%%%%%%%%%%%%%%%%%%%%%%%%%%%%%%%%%%%%%%%%%%%%%%%%
% ABBREVIATIONS
%%%%%%%%%%%%%%%%%%%%%%%%%%%%%%%%%%%%%%%%%%%%%%%%%%%%%%%%%%%%%%%%%

% Common abbreviations used in your paper
% Example: \newcommand{\ie}{i.e.,\xspace}
% Example: \newcommand{\eg}{e.g.,\xspace}
% Example: \newcommand{\cf}{cf.\xspace}
% Example: \newcommand{\wrt}{w.r.t.\xspace}

%%%%%%%%%%%%%%%%%%%%%%%%%%%%%%%%%%%%%%%%%%%%%%%%%%%%%%%%%%%%%%%%%
% DOMAIN-SPECIFIC TERMS
%%%%%%%%%%%%%%%%%%%%%%%%%%%%%%%%%%%%%%%%%%%%%%%%%%%%%%%%%%%%%%%%%

% Terms specific to your research domain
% Keep consistent formatting throughout the paper

%%%%%%%%%%%%%%%%%%%%%%%%%%%%%%%%%%%%%%%%%%%%%%%%%%%%%%%%%%%%%%%%%
% USAGE EXAMPLES
%%%%%%%%%%%%%%%%%%%%%%%%%%%%%%%%%%%%%%%%%%%%%%%%%%%%%%%%%%%%%%%%%
%
% GOOD:
%   \newcommand{\hypogenic}{\textsc{HypoGeniC}\xspace}
%   Usage: "We propose \hypogenic to generate hypotheses."
%   Result: "We propose HypoGeniC to generate hypotheses."
%
% GOOD:
%   \newcommand{\ours}{{\bf HypoGeniC}\xspace}
%   Usage: "Our method, \ours, outperforms baselines."
%   Result: "Our method, HypoGeniC, outperforms baselines."
%
% BAD (no \xspace):
%   \newcommand{\hypogenic}{\textsc{HypoGeniC}}
%   Usage: "\hypogenic outperforms..."
%   Result: "HypoGeniCoutperforms..." (missing space!)
%
%%%%%%%%%%%%%%%%%%%%%%%%%%%%%%%%%%%%%%%%%%%%%%%%%%%%%%%%%%%%%%%%%

%
% INSTRUCTIONS:
% 1. Use \textsc{} for method/dataset/model names (small caps)
% 2. Add \xspace at the end for automatic spacing
% 3. Group related terms together
% 4. Use consistent naming conventions

\usepackage{xspace}

%%%%%%%%%%%%%%%%%%%%%%%%%%%%%%%%%%%%%%%%%%%%%%%%%%%%%%%%%%%%%%%%%
% METHOD NAMES
%%%%%%%%%%%%%%%%%%%%%%%%%%%%%%%%%%%%%%%%%%%%%%%%%%%%%%%%%%%%%%%%%

% Your method name (replace with actual name)
\newcommand{\END}{\textsc{end}\xspace}          % Elicitation--Novelty Discriminator

% Discriminator quantities
\newcommand{\SG}{\textsc{sg}\xspace}            % Self-Generation utility
\newcommand{\LCG}{\textsc{lcg}\xspace}          % Length-Controlled Gain

% Task / setting names
\newcommand{\zeth}{\textsc{Zeth}\xspace}
\newcommand{\vocabzeth}{\textsc{vocab-zeth}\xspace}
\newcommand{\vocabes}{\textsc{vocab-es}\xspace}
\newcommand{\arithbten}{\textsc{arith-b10}\xspace}
\newcommand{\arithbnine}{\textsc{arith-b9}\xspace}
\newcommand{\arithbeleven}{\textsc{arith-b11}\xspace}

% Model names
\newcommand{\claude}{\textsc{Claude-Sonnet-4.5}\xspace}
\newcommand{\gpt}{\textsc{GPT-4.1}\xspace}
\newcommand{\qwen}{\textsc{Qwen2.5-7B-Instruct}\xspace}

% Abbreviations
\newcommand{\ie}{i.e.,\xspace}
\newcommand{\eg}{e.g.,\xspace}
\newcommand{\cf}{cf.\xspace}

%%%%%%%%%%%%%%%%%%%%%%%%%%%%%%%%%%%%%%%%%%%%%%%%%%%%%%%%%%%%%%%%%
% DATASET NAMES
%%%%%%%%%%%%%%%%%%%%%%%%%%%%%%%%%%%%%%%%%%%%%%%%%%%%%%%%%%%%%%%%%

% Dataset names (use \textsc for consistency)
% Example: \newcommand{\datasetone}{\textsc{Dataset-One}\xspace}
% Example: \newcommand{\datasettwo}{\textsc{Dataset-Two}\xspace}

%%%%%%%%%%%%%%%%%%%%%%%%%%%%%%%%%%%%%%%%%%%%%%%%%%%%%%%%%%%%%%%%%
% MODEL NAMES
%%%%%%%%%%%%%%%%%%%%%%%%%%%%%%%%%%%%%%%%%%%%%%%%%%%%%%%%%%%%%%%%%

% LLM and model names
% Example: \newcommand{\gpt}{\textsc{GPT-4o-mini}\xspace}
% Example: \newcommand{\llama}{\textsc{Llama-3.1-70B-Instruct}\xspace}
% Example: \newcommand{\claude}{\textsc{Claude-3.5-Sonnet}\xspace}

%%%%%%%%%%%%%%%%%%%%%%%%%%%%%%%%%%%%%%%%%%%%%%%%%%%%%%%%%%%%%%%%%
% BASELINE NAMES
%%%%%%%%%%%%%%%%%%%%%%%%%%%%%%%%%%%%%%%%%%%%%%%%%%%%%%%%%%%%%%%%%

% Baseline method names
% Example: \newcommand{\baseline}{\textsc{BaselineMethod}\xspace}
% Example: \newcommand{\fewshot}{\textsc{Few-shot}\xspace}

%%%%%%%%%%%%%%%%%%%%%%%%%%%%%%%%%%%%%%%%%%%%%%%%%%%%%%%%%%%%%%%%%
% ABBREVIATIONS
%%%%%%%%%%%%%%%%%%%%%%%%%%%%%%%%%%%%%%%%%%%%%%%%%%%%%%%%%%%%%%%%%

% Common abbreviations used in your paper
% Example: \newcommand{\ie}{i.e.,\xspace}
% Example: \newcommand{\eg}{e.g.,\xspace}
% Example: \newcommand{\cf}{cf.\xspace}
% Example: \newcommand{\wrt}{w.r.t.\xspace}

%%%%%%%%%%%%%%%%%%%%%%%%%%%%%%%%%%%%%%%%%%%%%%%%%%%%%%%%%%%%%%%%%
% DOMAIN-SPECIFIC TERMS
%%%%%%%%%%%%%%%%%%%%%%%%%%%%%%%%%%%%%%%%%%%%%%%%%%%%%%%%%%%%%%%%%

% Terms specific to your research domain
% Keep consistent formatting throughout the paper

%%%%%%%%%%%%%%%%%%%%%%%%%%%%%%%%%%%%%%%%%%%%%%%%%%%%%%%%%%%%%%%%%
% USAGE EXAMPLES
%%%%%%%%%%%%%%%%%%%%%%%%%%%%%%%%%%%%%%%%%%%%%%%%%%%%%%%%%%%%%%%%%
%
% GOOD:
%   \newcommand{\hypogenic}{\textsc{HypoGeniC}\xspace}
%   Usage: "We propose \hypogenic to generate hypotheses."
%   Result: "We propose HypoGeniC to generate hypotheses."
%
% GOOD:
%   \newcommand{\ours}{{\bf HypoGeniC}\xspace}
%   Usage: "Our method, \ours, outperforms baselines."
%   Result: "Our method, HypoGeniC, outperforms baselines."
%
% BAD (no \xspace):
%   \newcommand{\hypogenic}{\textsc{HypoGeniC}}
%   Usage: "\hypogenic outperforms..."
%   Result: "HypoGeniCoutperforms..." (missing space!)
%
%%%%%%%%%%%%%%%%%%%%%%%%%%%%%%%%%%%%%%%%%%%%%%%%%%%%%%%%%%%%%%%%%

%
% INSTRUCTIONS:
% 1. Use \textsc{} for method/dataset/model names (small caps)
% 2. Add \xspace at the end for automatic spacing
% 3. Group related terms together
% 4. Use consistent naming conventions

\usepackage{xspace}

%%%%%%%%%%%%%%%%%%%%%%%%%%%%%%%%%%%%%%%%%%%%%%%%%%%%%%%%%%%%%%%%%
% METHOD NAMES
%%%%%%%%%%%%%%%%%%%%%%%%%%%%%%%%%%%%%%%%%%%%%%%%%%%%%%%%%%%%%%%%%

% Your method name (replace with actual name)
\newcommand{\END}{\textsc{end}\xspace}          % Elicitation--Novelty Discriminator

% Discriminator quantities
\newcommand{\SG}{\textsc{sg}\xspace}            % Self-Generation utility
\newcommand{\LCG}{\textsc{lcg}\xspace}          % Length-Controlled Gain

% Task / setting names
\newcommand{\zeth}{\textsc{Zeth}\xspace}
\newcommand{\vocabzeth}{\textsc{vocab-zeth}\xspace}
\newcommand{\vocabes}{\textsc{vocab-es}\xspace}
\newcommand{\arithbten}{\textsc{arith-b10}\xspace}
\newcommand{\arithbnine}{\textsc{arith-b9}\xspace}
\newcommand{\arithbeleven}{\textsc{arith-b11}\xspace}

% Model names
\newcommand{\claude}{\textsc{Claude-Sonnet-4.5}\xspace}
\newcommand{\gpt}{\textsc{GPT-4.1}\xspace}
\newcommand{\qwen}{\textsc{Qwen2.5-7B-Instruct}\xspace}

% Abbreviations
\newcommand{\ie}{i.e.,\xspace}
\newcommand{\eg}{e.g.,\xspace}
\newcommand{\cf}{cf.\xspace}

%%%%%%%%%%%%%%%%%%%%%%%%%%%%%%%%%%%%%%%%%%%%%%%%%%%%%%%%%%%%%%%%%
% DATASET NAMES
%%%%%%%%%%%%%%%%%%%%%%%%%%%%%%%%%%%%%%%%%%%%%%%%%%%%%%%%%%%%%%%%%

% Dataset names (use \textsc for consistency)
% Example: \newcommand{\datasetone}{\textsc{Dataset-One}\xspace}
% Example: \newcommand{\datasettwo}{\textsc{Dataset-Two}\xspace}

%%%%%%%%%%%%%%%%%%%%%%%%%%%%%%%%%%%%%%%%%%%%%%%%%%%%%%%%%%%%%%%%%
% MODEL NAMES
%%%%%%%%%%%%%%%%%%%%%%%%%%%%%%%%%%%%%%%%%%%%%%%%%%%%%%%%%%%%%%%%%

% LLM and model names
% Example: \newcommand{\gpt}{\textsc{GPT-4o-mini}\xspace}
% Example: \newcommand{\llama}{\textsc{Llama-3.1-70B-Instruct}\xspace}
% Example: \newcommand{\claude}{\textsc{Claude-3.5-Sonnet}\xspace}

%%%%%%%%%%%%%%%%%%%%%%%%%%%%%%%%%%%%%%%%%%%%%%%%%%%%%%%%%%%%%%%%%
% BASELINE NAMES
%%%%%%%%%%%%%%%%%%%%%%%%%%%%%%%%%%%%%%%%%%%%%%%%%%%%%%%%%%%%%%%%%

% Baseline method names
% Example: \newcommand{\baseline}{\textsc{BaselineMethod}\xspace}
% Example: \newcommand{\fewshot}{\textsc{Few-shot}\xspace}

%%%%%%%%%%%%%%%%%%%%%%%%%%%%%%%%%%%%%%%%%%%%%%%%%%%%%%%%%%%%%%%%%
% ABBREVIATIONS
%%%%%%%%%%%%%%%%%%%%%%%%%%%%%%%%%%%%%%%%%%%%%%%%%%%%%%%%%%%%%%%%%

% Common abbreviations used in your paper
% Example: \newcommand{\ie}{i.e.,\xspace}
% Example: \newcommand{\eg}{e.g.,\xspace}
% Example: \newcommand{\cf}{cf.\xspace}
% Example: \newcommand{\wrt}{w.r.t.\xspace}

%%%%%%%%%%%%%%%%%%%%%%%%%%%%%%%%%%%%%%%%%%%%%%%%%%%%%%%%%%%%%%%%%
% DOMAIN-SPECIFIC TERMS
%%%%%%%%%%%%%%%%%%%%%%%%%%%%%%%%%%%%%%%%%%%%%%%%%%%%%%%%%%%%%%%%%

% Terms specific to your research domain
% Keep consistent formatting throughout the paper

%%%%%%%%%%%%%%%%%%%%%%%%%%%%%%%%%%%%%%%%%%%%%%%%%%%%%%%%%%%%%%%%%
% USAGE EXAMPLES
%%%%%%%%%%%%%%%%%%%%%%%%%%%%%%%%%%%%%%%%%%%%%%%%%%%%%%%%%%%%%%%%%
%
% GOOD:
%   \newcommand{\hypogenic}{\textsc{HypoGeniC}\xspace}
%   Usage: "We propose \hypogenic to generate hypotheses."
%   Result: "We propose HypoGeniC to generate hypotheses."
%
% GOOD:
%   \newcommand{\ours}{{\bf HypoGeniC}\xspace}
%   Usage: "Our method, \ours, outperforms baselines."
%   Result: "Our method, HypoGeniC, outperforms baselines."
%
% BAD (no \xspace):
%   \newcommand{\hypogenic}{\textsc{HypoGeniC}}
%   Usage: "\hypogenic outperforms..."
%   Result: "HypoGeniCoutperforms..." (missing space!)
%
%%%%%%%%%%%%%%%%%%%%%%%%%%%%%%%%%%%%%%%%%%%%%%%%%%%%%%%%%%%%%%%%%

%
% INSTRUCTIONS:
% 1. Use \textsc{} for method/dataset/model names (small caps)
% 2. Add \xspace at the end for automatic spacing
% 3. Group related terms together
% 4. Use consistent naming conventions

\usepackage{xspace}

%%%%%%%%%%%%%%%%%%%%%%%%%%%%%%%%%%%%%%%%%%%%%%%%%%%%%%%%%%%%%%%%%
% METHOD NAMES
%%%%%%%%%%%%%%%%%%%%%%%%%%%%%%%%%%%%%%%%%%%%%%%%%%%%%%%%%%%%%%%%%

% Your method name (replace with actual name)
\newcommand{\END}{\textsc{end}\xspace}          % Elicitation--Novelty Discriminator

% Discriminator quantities
\newcommand{\SG}{\textsc{sg}\xspace}            % Self-Generation utility
\newcommand{\LCG}{\textsc{lcg}\xspace}          % Length-Controlled Gain

% Task / setting names
\newcommand{\zeth}{\textsc{Zeth}\xspace}
\newcommand{\vocabzeth}{\textsc{vocab-zeth}\xspace}
\newcommand{\vocabes}{\textsc{vocab-es}\xspace}
\newcommand{\arithbten}{\textsc{arith-b10}\xspace}
\newcommand{\arithbnine}{\textsc{arith-b9}\xspace}
\newcommand{\arithbeleven}{\textsc{arith-b11}\xspace}

% Model names
\newcommand{\claude}{\textsc{Claude-Sonnet-4.5}\xspace}
\newcommand{\gpt}{\textsc{GPT-4.1}\xspace}
\newcommand{\qwen}{\textsc{Qwen2.5-7B-Instruct}\xspace}

% Abbreviations
\newcommand{\ie}{i.e.,\xspace}
\newcommand{\eg}{e.g.,\xspace}
\newcommand{\cf}{cf.\xspace}

%%%%%%%%%%%%%%%%%%%%%%%%%%%%%%%%%%%%%%%%%%%%%%%%%%%%%%%%%%%%%%%%%
% DATASET NAMES
%%%%%%%%%%%%%%%%%%%%%%%%%%%%%%%%%%%%%%%%%%%%%%%%%%%%%%%%%%%%%%%%%

% Dataset names (use \textsc for consistency)
% Example: \newcommand{\datasetone}{\textsc{Dataset-One}\xspace}
% Example: \newcommand{\datasettwo}{\textsc{Dataset-Two}\xspace}

%%%%%%%%%%%%%%%%%%%%%%%%%%%%%%%%%%%%%%%%%%%%%%%%%%%%%%%%%%%%%%%%%
% MODEL NAMES
%%%%%%%%%%%%%%%%%%%%%%%%%%%%%%%%%%%%%%%%%%%%%%%%%%%%%%%%%%%%%%%%%

% LLM and model names
% Example: \newcommand{\gpt}{\textsc{GPT-4o-mini}\xspace}
% Example: \newcommand{\llama}{\textsc{Llama-3.1-70B-Instruct}\xspace}
% Example: \newcommand{\claude}{\textsc{Claude-3.5-Sonnet}\xspace}

%%%%%%%%%%%%%%%%%%%%%%%%%%%%%%%%%%%%%%%%%%%%%%%%%%%%%%%%%%%%%%%%%
% BASELINE NAMES
%%%%%%%%%%%%%%%%%%%%%%%%%%%%%%%%%%%%%%%%%%%%%%%%%%%%%%%%%%%%%%%%%

% Baseline method names
% Example: \newcommand{\baseline}{\textsc{BaselineMethod}\xspace}
% Example: \newcommand{\fewshot}{\textsc{Few-shot}\xspace}

%%%%%%%%%%%%%%%%%%%%%%%%%%%%%%%%%%%%%%%%%%%%%%%%%%%%%%%%%%%%%%%%%
% ABBREVIATIONS
%%%%%%%%%%%%%%%%%%%%%%%%%%%%%%%%%%%%%%%%%%%%%%%%%%%%%%%%%%%%%%%%%

% Common abbreviations used in your paper
% Example: \newcommand{\ie}{i.e.,\xspace}
% Example: \newcommand{\eg}{e.g.,\xspace}
% Example: \newcommand{\cf}{cf.\xspace}
% Example: \newcommand{\wrt}{w.r.t.\xspace}

%%%%%%%%%%%%%%%%%%%%%%%%%%%%%%%%%%%%%%%%%%%%%%%%%%%%%%%%%%%%%%%%%
% DOMAIN-SPECIFIC TERMS
%%%%%%%%%%%%%%%%%%%%%%%%%%%%%%%%%%%%%%%%%%%%%%%%%%%%%%%%%%%%%%%%%

% Terms specific to your research domain
% Keep consistent formatting throughout the paper

%%%%%%%%%%%%%%%%%%%%%%%%%%%%%%%%%%%%%%%%%%%%%%%%%%%%%%%%%%%%%%%%%
% USAGE EXAMPLES
%%%%%%%%%%%%%%%%%%%%%%%%%%%%%%%%%%%%%%%%%%%%%%%%%%%%%%%%%%%%%%%%%
%
% GOOD:
%   \newcommand{\hypogenic}{\textsc{HypoGeniC}\xspace}
%   Usage: "We propose \hypogenic to generate hypotheses."
%   Result: "We propose HypoGeniC to generate hypotheses."
%
% GOOD:
%   \newcommand{\ours}{{\bf HypoGeniC}\xspace}
%   Usage: "Our method, \ours, outperforms baselines."
%   Result: "Our method, HypoGeniC, outperforms baselines."
%
% BAD (no \xspace):
%   \newcommand{\hypogenic}{\textsc{HypoGeniC}}
%   Usage: "\hypogenic outperforms..."
%   Result: "HypoGeniCoutperforms..." (missing space!)
%
%%%%%%%%%%%%%%%%%%%%%%%%%%%%%%%%%%%%%%%%%%%%%%%%%%%%%%%%%%%%%%%%%
